\documentclass[12pt]{article}
\usepackage[utf8]{inputenc}
\usepackage[T1]{fontenc}
\usepackage{hyperref}
\title{Totonaco del cerro Xinolat\'epetl Language Resources for K-12 Curriculum Development}
\author{Compiled by Claude Sonnet 5 (claude-sonnet-5) for Andr\'as Kornai}
\date{}
\begin{document}
\maketitle

\section{Introduction}
Totonaco del cerro Xinolat\'epetl (ISO 639-3 tqt; Glottocode west1511), also called Ozomatl\'an Totonac or Western Totonac, is a Totonacan language of Puebla, Mexico. Wikipedia cites about 1,600 speakers from the 2000 census. The searches turned up almost nothing beyond encyclopedia stubs, and no educational materials were found.

\section{Primary Educational Resources}
None found.

\section{Linguistic and Reference Materials}
\begin{itemize}
\item Ozomatl\'an Totonac (Wikipedia overview): \url{https://en.wikipedia.org/wiki/Ozomatl%C3%A1n_Totonac}
\end{itemize}

\section{Religious and Community Resources}
None found.

\section{Digital and Audio Resources}
None found.

\section{Curriculum Development Context}
This search found no school programmes, primers or dictionaries specific to this variety. Further work should look at Mexican bilingual education materials for Totonac (for example from the national indigenous-language agencies) and at documentation of neighbouring Totonac varieties. Those sources were not found here and are not cited.
\end{document}